%------------------------------------------------------------------------------%
\documentclass[a4paper,12pt,fleqn]{article}

\usepackage{calculo_varias_variaveis-1}

\ShowAnswers

%------------------------------------------------------------------------------%
\begin{document}

\makeHeader{CFVVI}{Prova 1 -- Turma 4}{05/05/2025}

% 01
%------------------------------------------------------------------------------%
\exercise{25}
Encontre uma parametrização para o movimento de uma partícula que começa no
ponto \((-2, 0)\) e traça a metade superior do círculo \(x^2+y^2=4\)
\clearpagequestiononly

\begin{answer}
  Metade superior do círculo de raio 2 centrado na origem

  Parametrização do círculo unitário centrado na origem
  \begin{align*}
    x      & = \cos(\theta) \\
    y      & = \sin(\theta) \\
    \theta & \in [0, 2\pi)
  \end{align*}
  Parametrização do círculo de raio 2 centrado na origem
  \begin{align*}
    x      & = 2\cos(\theta) \\
    y      & = 2\sin(\theta) \\
    \theta & \in [0, 2\pi)
  \end{align*}
  Parametrização da metade superior
  \begin{align*}
    x      & = 2\cos(\theta) \\
    y      & = 2\sin(\theta) \\
    \theta & \in [0, \pi]
  \end{align*}
  Essa parametrização começa em \((2, 0)\) e termina em \((-2, 0)\).
  Queremos inverter o sentido em que percorremos a trajetória
  \begin{align*}
    x      & = 2\cos(\pi -\theta) =    - 2 \cos(\theta)\\
    y      & = 2\sin(\pi -\theta) = \hpm 2 \sin(\theta)\\
    \theta & \in [0, \pi)
  \end{align*}
\end{answer}

% 02
%------------------------------------------------------------------------------%
\exercise{25}
Considerando a função
\(
  f(x, y) = \ln(\pi - x - y)
\)
\begin{enumerate}[label=(\alph*)]
  \item \ [10] Determine o domínio de $f$ e represente-o graficamente
  \item \ [5] Encontre a imagem de $f$
  \item \ [10] Caracterize todas as curvas de nível de $f$ e esboce três delas
\end{enumerate}
\begin{questiononly}
  \begin{multicols}{2}
    Domínio
    \vspace{-0.8\baselineskip}
    \begin{center}
      \includegraphics{axis_xy}
    \end{center}
    Curvas de nível
    \vspace{-0.8\baselineskip}
    \begin{center}
      \includegraphics{axis_xy}
    \end{center}
  \end{multicols}
  \clearpage
\end{questiononly}

\begin{answer}
  \begin{multicols}{2}
    Domínio
    \vspace{-0.8\baselineskip}
    \begin{center}
      \includegraphics{P1_T4_dominio}
    \end{center}
    Curvas de nível
    \vspace{-0.8\baselineskip}
    \begin{center}
      \includegraphics{P1_T4_curvas_nivel}
    \end{center}
  \end{multicols}

  \begin{enumerate}[label=(\alph*)]
    \item   O domínio de $f$ é o conjunto de pontos \((x, y)\) tais que
    \[
      \pi - x - y > 0
    \]
    \[
      y < \pi - x
    \]
    Portanto
    \[
      D_f = \{(x, y) \in \R^2 \st y < \pi - x\}
    \]

    \item A imagem de $f$ é o conjunto de valores que $f$ pode assumir.
    Como \( y < \pi - x \) assume valores no intervalo \((0, \infty)\)
    a imagem de $f$ é
    \(
      \text{Im}(f) = (-\infty, \infty)
    \)

    \item As curvas de nível são dadas por
    \[
      \ln(\pi - x - y) = c \qquad \text{ para } c \in \R
    \]
    \[
       \pi - x - y = e^c
    \]
    \[
      y = \pi -e^c - x
    \]
    isto é, retas com coeficiente angular $-1$ e intercepto $b = \pi - e^c$
    \begin{align*}
      b & = -2 & c & = \ln(\pi-(-2)) = \ln(\pi+2) & y & = -2 -x \\
      b & =  0 & c & = \ln(\pi)                   & y & = - x   \\
      b & =  2 & c & = \ln(\pi-2)                 & y & = 2 - x
    \end{align*}
  \end{enumerate}
\end{answer}

% 03
%------------------------------------------------------------------------------%
\exercise{30}
Calcule o limite ou mostre que ele não existe
\begin{enumerate}[label=(\alph*)]
  \item \ [15] \(\lim_{(x,y) \to (0, 0)} \frac{xy^2}{x^2 + y^4}\)
  \item \ [15] \(\lim_{(x,y) \to (4, 3)} \frac{\sqrt{x}-\sqrt{y+1}}{x-y-1}\)
\end{enumerate}
\clearpagequestiononly

\begin{answer}
\begin{enumerate}[label=(\alph*)]
\item Testando diferentes caminhos \\
  Caminho \(x = y^2\)
  \[
    \lim_{(x,y) \to (0,0)} \frac{xy^2}{x^2 + y^4} \evalat{x = y^2}{}
    = \lim_{y \to 0} \frac{y^2 y^2}{y^4 + y^4}
    = \lim_{y \to 0} \frac{y^4}{2y^4}
    = \lim_{y \to 0} \frac{1}{2}
    = \frac{1}{2}
  \]
  Caminho \(x = 0\)
  \[
    \lim_{(x,y) \to (0,0)} \frac{xy^2}{x^2 + y^4} \evalat{x = 0}{}
    = \lim_{y \to 0} \frac{0 \times y^2}{0 + y^4}
    = \lim_{y \to 0} 0
    = 0
  \]
  Como os limites são diferentes, o limite \textbf{não existe}

\item Multiplicando pelo conjugado
  \begin{align*}
    \frac{\sqrt{x}-\sqrt{y+1}}{x-y-1}
    & = \frac{\sqrt{x} - \sqrt{y+1}}{x - y - 1}
        \times
        \frac{\sqrt{x} + \sqrt{y+1}}{\sqrt{x} + \sqrt{y+1}} \\
    & = \frac{x - (y+1)}{(x - y - 1)\left(\sqrt{x} + \sqrt{y+1}\right)} \\
    & = \frac{1}{\sqrt{x} + \sqrt{y+1}}
  \end{align*}
  \begin{align*}
    \lim_{(x,y) \to (4, 3)} \frac{\sqrt{x}-\sqrt{y+1}}{x-y-1}
    & = \lim_{(x,y) \to (4, 3)} \frac{1}{\sqrt{x} + \sqrt{y+1}} \\
    & = \frac{1}{\sqrt{4} + \sqrt{3+1}} \\
    & = \frac{1}{2+2} = \frac{1}{4}
  \end{align*}
\end{enumerate}
\end{answer}

% 04
%------------------------------------------------------------------------------%
\exercise{20}
Encontre e esboce as curvas que representam o corte da superfície
\( x^2 + y^2 = 4 - z^2 \)
pelos planos
\(z=0\) e
\(y=0\)
\begin{questiononly}
  \begin{multicols}{2}
    \(z=0\)
    \vspace{-0.8\baselineskip}
    \begin{center}
      \includegraphics{axis_xy}
    \end{center}
    \(y=0\)
    \vspace{-0.8\baselineskip}
    \begin{center}
      \includegraphics{axis_xz}
    \end{center}
  \end{multicols}
\end{questiononly}

\begin{answer}
  \begin{multicols}{2}
    \(z=0\)
    \vspace{-0.8\baselineskip}
    \begin{center}
      \includegraphics{T4_conica_z0}
    \end{center}
    \(y=0\)
    \vspace{-0.8\baselineskip}
    \begin{center}
      \includegraphics{T4_conica_y0}
    \end{center}
  \end{multicols}

  No plano \(z=0\) a equação \( x^2 + y^2 = 4 - z^2 \) se reduz a
  \begin{align*}
    x^2 + y^2 & = 4 - 0 \\
    x^2 + y^2 & = 2^2
  \end{align*}
  que corresponde circunferência de raio 2 centrada na origem

  No plano \(y=0\) a equação \( x^2 + y^2 = 4 - z^2 \) se reduz a
  \begin{align*}
    x^2 + 0   & = 4 - z^2 \\
    x^2 + z^2 & = 2^2
  \end{align*}
  que corresponde circunferência de raio 2 centrada na origem
\end{answer}

%------------------------------------------------------------------------------%
\end{document}
%------------------------------------------------------------------------------%
